% !TeX root = ../papers/nc-p3.tex
\section{整六次点曲線の特異点と正規化の制限}
\label{sec:singular-sextic-normalization}
\begin{theorem}\label{thm:singular-sextic-normalization}
点ファイバーを持つ標準AS窓の点グラフが整な被約六次曲線 $C$ と交わらない有限部の和なら、
有限部は空で、$C$ は抽象曲線として整な特異平面四次である。
正規化の種数は0または1に限る。標準四移動が恒等でなければ、標準平面上の座標を取ると
$[u:v:w]\mapsto[\eta u:v:w]$、$\eta$ の位数は2または4であり、$K^2=1$ である。
位数4の特異模型は $u^4+F_4(v,w)=0$ で、二項四次の根の重複型は $(2,1,1)$ または $(3,1)$ である。
\end{theorem}
\begin{proof}[三次元接点の高さと極]
平面六次は隣接平面積の次数で除かれる。接次元3の点があれば
命題\ref{prop:automatic-spatial-tangent-mark}のOre表示を取る。
その水平射影の像は整な非退化平面曲線 $D$ である。三生成基底が平面グラフなら
定理\ref{thm:integral-septic-boundary}と同じく曲面が入る。残る基底は三次曲線のグラフで、
二次曲線なら定理\ref{thm:singular-integral-quintic-exclusion}の二次の極上限により
持ち上がる曲線の次数は高々4となる。実際、二つの高さは一般点で水平点から一意に戻るため
水平点曲線の関数体は $C$ の関数体である。従って $D$ は整三次で、基底の移動を $g$ とする。

$g$ の固定点 $x$ では $v(x)\wedge d_i(v(x),v(gx))\ne0$ である。零なら混合式が解を持ち、
両方の高さを同時に変える直線が点グラフの成分になるからである。
一般点で $a\,v(gx)-b\,v(x)=d_i(v(x),v(gx))$ を満たすので、
一つが単元となる成分で割ると、固定点の局所イデアルは $1/a$ で生成される（局所自明化を用いる）。
その正規化への引戻しは正確に高さの極因子となる。従って元の曲線の次数は
$3$ とこの極次数の和である。高さと水平三座標で基点のない射を与え、一般点では逆もあるからである。

滑らかな三次の基底なら定理\ref{thm:all-elliptic-ore-windows}を使える。
固定点のない移動では極はなく、固定点がある位数3,4の基底では補正が零に同値で、
各固定点に直線が生じる。残る反転は四つの単純固定点を持ち次数は7である。
節点三次の自己同型で二枝を保つ非恒等なものは正規化で $z\mapsto cz$、
節点での固定イデアルは二枝の極大イデアルで主でない。枝を交換すれば $z\mapsto c/z$ となり、
枝での固定イデアルの引戻しは各長さ1、その外に単純固定点二つで極次数は4である。
尖点三次の滑らかな部分で移動は $z\mapsto az+b$ と書ける。
三生成ASの偏極条件で $g^*L/L$ は非自明かつ $g^*$ 不変な次数零類である
（補題\ref{lem:ternary-finite-graph-test}の三次関係の行列式、または補題\ref{lem:reduced-ternary-fourth-shift}の束計算）。
$\operatorname{Pic}^0=\mathbb G_a$ への作用は $a$ または $a^{-1}$ 倍である
（尖点での正規化の一次jetの単元商）。よって $a=1$, $b\ne0$ である。
尖点のパラメータ $t=1/z$ に対し $t\mapsto t/(1+bt)$ で、
$t^2,t^3$ の差の位数は3,4となる。$k[[t^2,t^3]]$ に位数1の元がなく、この固定イデアルも主でない。
以上で六次では接次元3の点はない。

\smallskip
\noindent\textit{標準平面と固定直線。}
従って全特異点は局所平面、グラフは局所完全交叉である。
超過交叉の長さ $2p_a-6$ とCastelnuovo上界から $p_a=3,4$ であり、
種数4なら全埋込みが標準埋込みで同じ二次曲面の対角を含むという既出の矛盾になる。
種数3では有限部は空である。標準束は基点自由である。実際任意の点 $x$ について
$\mathcal Hom(\mathfrak m_x,\mathcal O_C)$ の大域切断は定数のみである。
滑らかな点で単純極一つを持つ函数は次数1の射を与えてしまい、特異点では
$(\mathcal O:\mathfrak m_x)=\operatorname{End}(\mathfrak m_x)$ である
（像に単元があれば極大イデアルが主になる）。後者は有限な整な部分正規化の函数層である。
双対性から基点自由性を得る。三切断による次数4の標準写像は二次曲線への有限平坦二重被覆か、
算術種数が等しいため整平面四次への同型である。二重被覆では定理\ref{thm:smooth-sextic-exclusion}の
鉛筆分裂の証明がそのまま成立する。$L_i/\omega_C$ の有効な長さ2のCartier因子は埋込みに反し、
有限平坦な二重被覆のtrace分解から $\omega_C=J^2$ である。
以後は直線上の束の分裂と各長さ2の環上の乗法しか用いず、三次の次元上限18に達する。

標準平面で $\tau=\sigma^{-1}$ が非恒等なら、固定切断は $C$ 上で長さ4のCartier因子を切る。
射影変換が固定直線を持たなければ、平面での固定スキームの長さは3で足りない
（接束の切断の孤立零の最高Chern数、またはJordan形での小行列式）。
固定直線を持つ非半単純な変換は一つの推移 $u\mapsto u+v$ なので、
不変な四次曲線の式は $v,w$ だけになり整でない。
半単純なら固定直線の四交点で全長が尽き、中心点は曲線の外である。
$u^4$ と $u$ を含まない部分がともに非零だから倍率の四乗は1である。
固定因子は標準直線切断なので定理\ref{thm:fixed-differential}より $K^2=1$。
位数4の式で全根の多重度が偶数なら可約で、全て単純なら滑らかである。

\smallskip
\noindent\textit{正規化の種数2の排除。}
この場合はただ一つの節点か尖点がある。定理\ref{thm:fixed-differential}の恒等の場合の式で、
$\deg M\ge 2-1=1$ かつ $\deg M(T)=\deg\omega_C=4$ なので $\tau\ne1$ である。
節点の二枝を保つ有限位数の非恒等写像は両枝で接倍率が1でなく、
固定イデアルが主にならない。交換する場合は接作用の固有値に1が必要なので、
その平方の枝接倍率が1となり、有限位数性から $\tau^2=1$ である。
尖点に位数4の作用があれば、正規化の一様化元を固有元にして、
$t^2,t^3$ の差がともに非零倍率となるため主でない。従ってどちらも対合である。
節点では固定長2と二つの滑らかな固定点、尖点では固定長3と一つの滑らかな固定点を得る。
正規化 $N$ 上の対合はいずれも二つの固定点を持ち、Hurwitzにより商の種数は1である。

一般の滑らかな $p$ で定理\ref{thm:nonhyperelliptic-sextic-exclusion}の
長さ2の反対階数低下因子 $D_p$ と三割線を取る。ここで要る特異曲線への移行を説明する。
どの六次埋込みにも二次式はない。滑らかな二次曲面上の種数3の型は $(2,4)$ またはその逆で二重被覆になる。
錐では次数6の曲線のWeil類は三超平面の和なので算術種数4である。
$T=L_2/\omega_C$ と $\omega_C/T$ は非有効な次数2の束である。
曲線の二次式がなく、全二次切断が十次元なのでHilbert--Burchの四三次小行列による記述は同じである。
線型三行四列は曲線上で常に階数2で、階数1なら接次元が3になってしまう。
行列式三次 $S_p$ の非消滅について、中支持4の証明は次数や滑らかさによらず、
中支持3のときに平面グラフから得る二点 $a,b$ は一般の $p$ の点列上にあり滑らかである。
同じ共通三文字を用いる次数7の束は非常に豊富である（長さ2のイデアルでひねった後の
$H^1$ の双対は負次数の階数1無捩れ層の切断となる）。従って三二次式の次数8との矛盾も同じである。
$D_p$ が一般に特異点を避け、$\mathcal O(D_p)=\mathcal O(p+\tau p)/K$ となる理由は
定理\ref{thm:singular-integral-quintic-exclusion}の二核束の議論による。

三割線との長さ3の交叉のうち $D_p$ を除いた長さ1の残余は一般には滑らかな可動点 $r(p)$ である。
もしある固定点 $z$ を通る直線ばかりなら、$z$ からの射影の一般ファイバーが
長さ2の $D_p$ を含む。射影平面を張る像に対し射影束の次数は $6-\operatorname{mult}_z C\le5$ なので、
分離性から二次曲線への二重被覆しかない。そのとき曲線が二次錐に入ってしまう。
$D_p$ の一般台は可動であるから上のファイバーは一般に $z$ 以外の滑らかな部分で数えてよい。
この議論で特異点が残余に常駐することも除かれる。
従って一般の $p$ で交叉は有効Cartier因子となる。たとえ余点が $D_p$ に重なっても
滑らかな点では長さの加算が因子の加算となる。直線を含む平面とRiemann--Rochから
$D_p+r(p)\sim T+s(p)$、$s(p)$ は有効な次数1の束が定める滑らかな点である。
一意な残余と一意な有効切断により両点は関数体上定義され、$N$ の自己写像に延びる。
$|T+s|$ は任意の滑らかな $s$ で一点なので $s(p)$ も非定数である。
種数2ではHurwitzにより $r,s$ はいずれも自己同型である。

$g=rs^{-1}$ とすると、$N$ のJacobianに
\[
                  (1-g_*)s_*=1+\tau_*
\]
が成立する。これは一般 $p$ での $p+\tau p+r(p)-s(p)\sim T K$ を正規化へ引き戻したもの。
$T K$ は非有効である。そうでなければ非超楕円性からその有効因子 $Z$ は唯一であり、
一般の $s$ では $|T K+s|=\{Z+s\}$ となる（双対な次数2の束の切断も高々一つ）。
これは可動な相異なる二点 $p,\tau p$ を含めない。従って $s$ は恒等でも $\tau$ でもない。
そうであれば上の次数3の因子から $s(p)$ を引いて $T K$ が有効になるためである。

標数零の有限個のデータを複素数へ埋め、有限自己同型群に不変なHermite内積を取る。
Jacobian等式に随伴を右から掛けると $g+g^{-1}=-2\tau$ である。
$\tau$ の $+1,-1$ 固有空間が各一次元のため、$g=-\tau$、$s$ は
この順に $1,\xi$ で作用する。有限位数の非恒等写像 $s,\tau s$ の固定点数は
Lefschetz数 $-2\operatorname{Re}\xi,2\operatorname{Re}\xi$ に等しい
（両方とも全固定点が単純）。従って両者とも固定点がなく、$\xi=\pm\mathrm i$。
種数2のJacobian作用は忠実なので $s^2=\tau$ である。
巡回四重被覆 $N\to N/\langle s\rangle$ の唯一の分岐軌道は $\tau$ の二固定点で、
分岐指数2である。これは不可能である。実際Kummer表示の四乗される根号の中の有理函数は
その一箇所でのみ位数が $2\bmod4$、他はすべて $0\bmod4$ となり主因子の次数が零にならない。
滑らかな種数3の場合は既出の全排除を用いれば、正規化種数は0,1のみである。
\end{proof}
